% 数学ノート集・共通テンプレート v2（積分論2026年度第1部の版面を参考）
% A5判。章内で定義・定理・例・注意等の番号を共通化。
% コンパイル： platex -interaction=nonstopmode math_note_template_v2.tex (2回)
%              dvipdfmx math_note_template_v2.dvi
\documentclass[a5paper,10pt,onecolumn]{jsarticle}
\usepackage[dvipdfmx,top=22mm,bottom=20mm,inner=19mm,outer=18mm,headheight=13pt,headsep=6mm]{geometry}
\usepackage{amsmath,amssymb,amsthm,mathtools}
\usepackage{enumitem,fancyhdr,titlesec}
\usepackage[dvipdfmx,bookmarks=true,colorlinks=false]{hyperref}
\usepackage{pxjahyper}

% ----- 差し替える文書情報 -----
\newcommand{\NoteTitle}{ガロア理論と方程式の可解性}
\newcommand{\NoteSeries}{数学ノート集}

% ----- 文字・組版 -----
% 参考PDFの質素な明朝＋Computer Modernの印象を引き継ぎながら，
% A5判で窮屈にならないよう本文の行間を少し確保する．
\linespread{1.13}
\setlength{\parindent}{1zw}
\setlength{\parskip}{0pt}
\setlength{\abovedisplayskip}{6pt plus 2pt minus 2pt}
\setlength{\belowdisplayskip}{6pt plus 2pt minus 2pt}
\setlength{\abovedisplayshortskip}{3pt plus 2pt minus 1pt}
\setlength{\belowdisplayshortskip}{4pt plus 2pt minus 1pt}
\setlist[enumerate]{topsep=3pt,itemsep=1pt,parsep=0pt,leftmargin=2.4zw}
\titleformat{\section}{\normalfont\normalsize\bfseries}{\thesection}{0.8zw}{}
\titlespacing*{\section}{0pt}{2.0ex plus .4ex}{.9ex}
\titleformat{\subsection}{\normalfont\normalsize\bfseries}{\thesubsection}{0.7zw}{}
\titlespacing*{\subsection}{0pt}{1.5ex plus .3ex}{.6ex}
\raggedbottom

% ----- 柱とページ番号：ページ上部、罫線なし -----
\fancypagestyle{mathnote}{%
  \fancyhf{}%
  \fancyhead[L]{\footnotesize \NoteSeries\ （\NoteTitle）}%
  \fancyhead[R]{\footnotesize \thepage}%
  \renewcommand{\headrulewidth}{0pt}%
  \renewcommand{\footrulewidth}{0pt}%
}
\pagestyle{mathnote}

% ----- 節内通し番号の定義・定理群 -----
\newtheoremstyle{mathnotetheorem}
 {4.5pt}{4.5pt}{\normalfont}{}%
 {\bfseries}{}{0.65zw}%
 {\thmname{#1}\thmnumber{ #2}\thmnote{（#3）}}
\theoremstyle{mathnotetheorem}
\newtheorem{theorem}{定理}[section]
\newtheorem{definition}[theorem]{定義}
\newtheorem{lemma}[theorem]{補題}
\newtheorem{proposition}[theorem]{命題}
\newtheorem{corollary}[theorem]{系}
\newtheorem{example}[theorem]{例}
\newtheorem{remark}[theorem]{注意}
\renewcommand{\proofname}{証明}
\renewcommand{\qedsymbol}{$\square$}
\makeatletter
\renewenvironment{proof}[1][\proofname]{%
  \par\pushQED{\qed}\normalfont\topsep4\p@\@plus2\p@\relax
  \trivlist\item[\hskip\labelsep\bfseries #1]\ignorespaces
}{\popQED\endtrivlist\@endpefalse}
\makeatother

% ----- 使うノートで適宜拡張する記号 -----
\newcommand{\R}{\mathbb{R}}
\newcommand{\norm}[1]{\lVert #1\rVert}
\newcommand{\abs}[1]{\lvert #1\rvert}

\begin{document}
\section*{読了後にできること}
本稿を読むための前提は，Galois対応，有限群の可解性．
本稿の具体的な到達課題は，具体的な多項式のGalois群を決め，根号可解性を判定する．

\section{根号による可解性}

本稿では全ての体は標数 $0$ とする．
有限次Galois拡大，分解体，Galois対応および
有限群の正規部分群の基本理論を前提とする．

\begin{definition}
体 $K$ 上の多項式 $f(X)\in K[X]$ が
\emph{根号によって可解}であるとは，
ある体拡大の塔
\[
 K=K_0\subset K_1\subset\cdots\subset K_r
\]
が存在し，各 $i$ について
\[
 K_i=K_{i-1}(\alpha_i),\qquad
 \alpha_i^{n_i}\in K_{i-1}
 \quad(n_i\ge2)
\]
を満たし，$f$ の全ての根が $K_r$ に属することをいう．
このような塔を\emph{根号拡大の塔}という．
\end{definition}

\begin{remark}
必要なら $1$ の冪根を付け加える操作も
$\zeta_n^n=1$ による根号拡大である．
従って定義では，四則演算と有限回の根号抽出，
および根号を取る際の異なる値の選択を許している．
\end{remark}

\begin{definition}
有限群 $G$ の導来列を
\[
 G^{(0)}=G,\qquad
 G^{(i+1)}=[G^{(i)},G^{(i)}]
\]
と定める．
ある $r$ について $G^{(r)}=\{1\}$ となるとき，
$G$ を\emph{可解群}という．
\end{definition}

\begin{lemma}[可解群の基本性質]\label{lem:solvable}
可解群の部分群および商群は可解である．
また $N\triangleleft G$ について
$N$ と $G/N$ がともに可解なら $G$ も可解である．
有限個の可解群の直積も可解である．
\end{lemma}
\begin{proof}
$H\le G$ なら
$H^{(i)}\subset G^{(i)}$ が帰納的に従うので，
$G$ が可解なら $H$ も可解である．
全射準同型 $\pi:G\to Q$ に対して
$\pi(G^{(i)})=Q^{(i)}$ だから，
$G$ が可解なら $Q$ も可解である．

$N\triangleleft G$ とし，$G/N$ が可解なら
ある $r$ で $(G/N)^{(r)}=1$ となる．
従って $G^{(r)}\subset N$ である．
さらに $N$ が可解で $N^{(s)}=1$ なら，
$G^{(r+s)}=(G^{(r)})^{(s)}
 \subset N^{(s)}=1$ となる．
最後に直積について
$(G_1\times\cdots\times G_m)^{(i)}
 =G_1^{(i)}\times\cdots\times G_m^{(i)}$
を用いればよい．
\end{proof}

\begin{lemma}[可解群と巡回素数位数の正規列]\label{lem:composition}
有限群 $G$ が可解であることと，
部分群の列
\[
 G=G_0\triangleright G_1\triangleright\cdots
 \triangleright G_r=\{1\}
\]
で各 $G_i/G_{i+1}$ が素数位数の巡回群であるものが
存在することは同値である．
ここで $G_{i+1}$ は $G_i$ の正規部分群であればよく，
$G$ 全体の正規部分群である必要はない．
\end{lemma}
\begin{proof}
列が存在すれば，各商群はAbel群なので可解である．
補題\ref{lem:solvable}の拡大に関する性質を
帰納的に適用すれば $G$ は可解となる．
逆に $G$ を有限可解群とする．
有限群には組成列が存在する．
その各組成因子は $G$ の部分群の商群であるから，
補題\ref{lem:solvable}によって可解である．
一方，単純な非自明可解群 $S$ について
$[S,S]$ は $S$ の正規部分群なので
$\{1\}$ または $S$ に等しい．
後者なら導来列は零にならず可解性に反するため
$[S,S]=1$ であり，$S$ はAbel群である．
有限単純Abel群は素数位数の巡回群だから，
各組成因子は巡回素数位数である．
従って組成列が所望の列となる．
\end{proof}

\section{1の冪根を含む巡回拡大}

\begin{lemma}[根号拡大のGalois群]\label{lem:kummer-abelian}
体 $F$ が全ての $n$ 乗根の1を含むとする．
$a_1,\dots,a_s\in F^\times$ に対して
\[
 E=F(\alpha_1,\dots,\alpha_s),\qquad
 \alpha_j^n=a_j
\]
と置く．このとき $E/F$ は有限Galois拡大であり，
$\operatorname{Gal}(E/F)$ はAbel群である．
\end{lemma}
\begin{proof}
各多項式 $X^n-a_j$ は標数 $0$ なので分離的である．
その根は $\alpha_j\zeta$（$\zeta^n=1$）であり，
仮定から全て $E$ に属する．
従って $E$ はこれらの多項式の分解体なので
有限Galois拡大である．
任意の $\sigma\in\operatorname{Gal}(E/F)$ に対して
$\sigma(\alpha_j)^n=\sigma(a_j)=a_j$ なので
$\sigma(\alpha_j)=\zeta_{\sigma,j}\alpha_j$
となる $n$ 乗根 $\zeta_{\sigma,j}\in F$ が存在する．
写像
\[
 \operatorname{Gal}(E/F)\to(\mu_n)^s,\qquad
 \sigma\mapsto(\zeta_{\sigma,1},\dots,\zeta_{\sigma,s})
\]
は準同型である．
実際 $\zeta_{\tau,j}\in F$ なので
$(\sigma\tau)(\alpha_j)
 =\sigma(\zeta_{\tau,j}\alpha_j)
 =\zeta_{\tau,j}\zeta_{\sigma,j}\alpha_j$
となる．
全ての $\alpha_j$ が固定されれば $\sigma$ は
$E$ 全体を固定するのでこの写像は単射である．
右辺はAbel群だから左辺もAbel群である．
\end{proof}

\begin{lemma}[巡回素数次拡大は根号拡大]\label{lem:cyclic-radical}
$E/F$ を素数 $\ell$ 次の巡回Galois拡大とし，
$F$ が原始 $\ell$ 乗根 $\zeta_\ell$ を含むとする．
このときある $\alpha\in E$ が存在して
\[
 E=F(\alpha),\qquad \alpha^\ell\in F.
\]
\end{lemma}
\begin{proof}
$\sigma$ を $\operatorname{Gal}(E/F)$ の生成元とする．
$E$ は $F$ 上の有限次元線形空間であり，
$\sigma$ は $F$ 線形写像で $\sigma^\ell=I$ を満たす．
$X^\ell-1$ は標数0で重根を持たず，
その根は全て $F$ に属するので
$\sigma$ は $F$ 上対角化可能である．
$\sigma\ne I$ だから，固有値1以外の固有値
$\zeta_\ell^j$ を持つ非零固有ベクトル
$\alpha\in E$ が存在する．
$\ell$ は素数だから $j\not\equiv0\pmod\ell$ であり，
$\zeta_\ell^j$ は原始 $\ell$ 乗根である．
\[
 \sigma(\alpha^\ell)
 =(\zeta_\ell^j)^\ell\alpha^\ell=\alpha^\ell
\]
なので $\alpha^\ell\in F$ となる．
また $\alpha\notin F$ であり，
$[E:F]=\ell$ は素数なので
$F(\alpha)=E$ である．
\end{proof}

\begin{lemma}[1の冪根の添加]\label{lem:cyclotomic}
任意の正整数 $N$ について，
$K(\mu_N)/K$ は有限Abel型Galois拡大である．
\end{lemma}
\begin{proof}
$X^N-1$ は標数0で分離的であり，
$K(\mu_N)$ はその分解体なので有限Galois拡大である．
任意の自己同型は原始 $N$ 乗根 $\zeta_N$ を
別の原始 $N$ 乗根へ写すので，
\[
 \operatorname{Gal}(K(\mu_N)/K)
 \hookrightarrow(\mathbb Z/N\mathbb Z)^\times
\]
という単射準同型が得られる．
右辺はAbel群なので左辺もAbel群である．
\end{proof}

\section{根号可解性からGalois群の可解性へ}

\begin{theorem}[必要性]\label{thm:necessary}
$K$ を標数0の体，$f\in K[X]$ を分離的多項式とし，
その分解体を $L$ とする．
$f$ が根号によって可解なら，
$\operatorname{Gal}(L/K)$ は可解群である．
\end{theorem}
\begin{proof}
根号拡大の塔
\[
 K=E_0\subset E_1\subset\cdots\subset E_r,\qquad
 E_i=E_{i-1}(\alpha_i),\quad \alpha_i^{n_i}\in E_{i-1}
\]
で $L\subset E_r$ となるものを取る．
$\alpha_i=0$ の段階は $E_i=E_{i-1}$ であって
拡大を生じないので，そのような段階を除いてよい．
従って以下では全ての $\alpha_i$ と
$a_i=\alpha_i^{n_i}$ は非零と仮定する．
$N$ を全ての $n_i$ の公倍数とし，
$M_0=K(\mu_N)$ と置く．
補題\ref{lem:cyclotomic}によって
$M_0/K$ は有限GaloisでGalois群は可解である．

帰納的に，$M_{i-1}/K$ は有限Galoisで可解な
Galois群を持ち，$E_{i-1}\subset M_{i-1}$ を満たすように
$M_{i-1}$ を構成できたとする．
$a_i=\alpha_i^{n_i}\in E_{i-1}$ と置く．
有限集合
\[
 A_i=\{\sigma(a_i):\sigma\in\operatorname{Gal}(M_{i-1}/K)\}
\]
を考える．$M_i$ を $M_{i-1}$ 上の多項式
\[
 P_i(X)=\prod_{a\in A_i}(X^{n_i}-a)
\]
の分解体とする．
$M_{i-1}$ は $\mu_{n_i}$ を含むので，
補題\ref{lem:kummer-abelian}から
$\operatorname{Gal}(M_i/M_{i-1})$ はAbel群である．
また $\operatorname{Gal}(M_{i-1}/K)$ は
$A_i$ を置換するため，
$P_i(X)$ の係数は $K$ に固定される．
従って $M_i$ は $K$ 上の有限個の分離的多項式を
全て分解する体であり，$M_i/K$ は有限Galoisである．
Galois群の完全列
\[
 1\to\operatorname{Gal}(M_i/M_{i-1})
 \to\operatorname{Gal}(M_i/K)
 \to\operatorname{Gal}(M_{i-1}/K)\to1
\]
と補題\ref{lem:solvable}によって，
$\operatorname{Gal}(M_i/K)$ は可解である．
$a_i\in A_i$ なので $M_i$ は $\alpha_i$ を含み，
$E_i\subset M_i$ となる．
これで帰納法が成立する．

最後に $L\subset E_r\subset M_r$ であり，
$L/K$ はGaloisだから，
制限写像
$\operatorname{Gal}(M_r/K)\to\operatorname{Gal}(L/K)$
は全射である．
可解群の商群は可解なので
$\operatorname{Gal}(L/K)$ は可解である．
\end{proof}

\begin{remark}
証明では，根号拡大の塔そのものが
Galois拡大の塔であるとは仮定していない．
各段階で共役な根号を全て付け加えて
$K$ 上正規な拡大に作り直すことが重要である．
\end{remark}

\section{Galois群の可解性から根号可解性へ}

\begin{theorem}[十分性]\label{thm:sufficient}
$K$ を標数0の体，$f\in K[X]$ を分離的多項式とし，
その分解体を $L$ とする．
$\operatorname{Gal}(L/K)$ が可解群なら，
$f$ は根号によって可解である．
\end{theorem}
\begin{proof}
$G=\operatorname{Gal}(L/K)$ とし，
$N=|G|$ と置く．$N=1$ の場合は $L=K$ なので明らかである．
$M=K(\mu_N)$ とすると
$M/K$ はAbel型Galois拡大であり，
$M$ は $1$ の $\ell$ 乗根を
全て含む（$\ell$ が $N$ の素因数である場合）．
$LM/M$ はGaloisであり，
$\operatorname{Gal}(LM/M)$ は
$G$ の部分群に同型なので可解である．
補題\ref{lem:composition}によって
\[
 G_0\triangleright G_1\triangleright\cdots
 \triangleright G_s=1,\qquad
 G_i/G_{i+1}\cong\mathbb Z/\ell_i\mathbb Z
\]
という正規列が得られる．
ここで各 $\ell_i$ は $N$ の素因数である．
Galois対応によって
\[
 M=F_0\subset F_1\subset\cdots\subset F_s=LM
\]
という体の塔が得られ，
各 $F_{i+1}/F_i$ は素数 $\ell_i$ 次の巡回Galois拡大である．
全ての $F_i$ は $\mu_{\ell_i}$ を含むので，
補題\ref{lem:cyclic-radical}から
\[
 F_{i+1}=F_i(\alpha_i),\qquad
 \alpha_i^{\ell_i}\in F_i
\]
となる．
従って $LM/M$ は根号拡大の塔である．
また $M=K(\mu_N)$ 自身も
$K$ に $N$ 乗根 $\zeta_N$（$\zeta_N^N=1$）
を付け加える根号拡大である．
よって $K$ から $LM$ への塔全体は根号拡大の塔となる．
$L\subset LM$ だから，$f$ の全ての根はこの塔に属し，
$f$ は根号によって可解である．
\end{proof}

\begin{theorem}[根号可解性のGalois理論的判定]\label{thm:iff}
標数0の体 $K$ 上の分離的多項式 $f$ の分解体を $L$ とする．
このとき次の条件は同値である．
\begin{enumerate}
\item $f$ は根号によって可解である．
\item $\operatorname{Gal}(L/K)$ は可解群である．
\end{enumerate}
\end{theorem}
\begin{proof}
(1) から (2) は定理\ref{thm:necessary}，
(2) から (1) は定理\ref{thm:sufficient}による．
\end{proof}

\section{対称群と交代群}

\begin{lemma}[交代群の基本性質]\label{lem:alternating}
$n\ge2$ のとき
$A_n=\ker(\operatorname{sgn}:S_n\to\{\pm1\})$
は $S_n$ の指数2の正規部分群である．
また $A_n$ は3-cycle全体によって生成される．
\end{lemma}
\begin{proof}
符号写像は全射準同型なので，
その核は指数2の正規部分群である．
任意の偶置換は偶数個の互換の積で書ける．
二つの互換の積は，
互換が同じなら恒等置換，
一点を共有するなら3-cycle，
互いに交わらないなら
\[
 (ab)(cd)=(ab)(bc)(bc)(cd)
\]
のように3-cycle二つの積に書き直せる．
実際，例えば
$(ab)(cd)=(acb)(acd)$
という直接計算で確認できる表示を用いてもよい．
従って偶置換は3-cycleの積で表され，
$A_n$ は3-cycle全体で生成される．
\end{proof}

\begin{lemma}[ $A_5$ の共役類]\label{lem:a5-classes}
$A_5$ の共役類の大きさは
\[
 1,\quad15,\quad20,\quad12,\quad12
\]
である．それぞれ恒等置換，二つの互換の積，
3-cycle，および二種類の5-cycleに対応する．
\end{lemma}
\begin{proof}
$|A_5|=5!/2=60$ である．
偶置換の可能な巡回型は
恒等置換，
$(ab)(cd)$ 型，
3-cycle型，
5-cycle型の四種類である．
それぞれの元の個数は
\[
 1,\quad\frac12\binom52\binom32=15,\quad
 \binom53\cdot2=20,\quad(5-1)!=24.
\]
3-cycle $(123)$ の $S_5$ における中心化群は
$\langle(123)\rangle\times\langle(45)\rangle$ で位数6である．
その偶置換部分は位数3なので，
$A_5$ における共役類の大きさは $60/3=20$ である．
二重互換 $(12)(34)$ の $S_5$ における中心化群は
位数8で，偶置換部分は位数4である．
従って $A_5$ における共役類の大きさは $60/4=15$ である．
5-cycleの中心化群はその生成する位数5の巡回群であり，
全て偶置換なので，$A_5$ における共役類の大きさは
$60/5=12$ である．
5-cycleは24個あるから二つの共役類に分かれる．
\end{proof}

\begin{theorem}[ $A_5$ の単純性]\label{thm:a5-simple}
$A_5$ は非可換単純群である．
従って $A_5$ は可解群ではない．
\end{theorem}
\begin{proof}
$N\triangleleft A_5$ を正規部分群とする．
$N$ は共役類の和集合であり単位元を含むから，
補題\ref{lem:a5-classes}によってその位数は
\[
 |N|=1+15\varepsilon_1+20\varepsilon_2
       +12\varepsilon_3+12\varepsilon_4,
 \qquad \varepsilon_i\in\{0,1\}
\]
の形である．
Lagrangeの定理により $|N|$ は $60$ の約数である．
しかし上の16通りの和のうち，
$60$ の約数となるものは $1$ と $60$ のみである．
実際，真の非自明な和は少なくとも $13$ であり，
$60$ の真の約数で $13$ 以上のものは
$15,20,30$ だけであるが，
上の和はそのいずれにも等しくならない．
従って $N=\{1\}$ または $N=A_5$ であり，
$A_5$ は単純である．
また $(123)$ と $(345)$ は可換でないので
$A_5$ は非可換である．
非可換単純群の交換子部分群は
非自明な正規部分群なので群全体に等しく，
導来列は零にならない．
従って $A_5$ は可解でない．
\end{proof}

\begin{corollary}\label{cor:sn-nonsolvable}
$n\ge5$ に対して $S_n$ は可解群ではない．
\end{corollary}
\begin{proof}
$S_n$ は最初の5個の文字だけを置換する
部分群として $S_5$ を含み，$S_5$ は $A_5$ を含む．
もし $S_n$ が可解なら
補題\ref{lem:solvable}からその部分群 $A_5$ も可解となるが，
定理\ref{thm:a5-simple}に矛盾する．
\end{proof}

\section{一般の方程式のGalois群}

\begin{theorem}[一般の $n$ 次多項式のGalois群]\label{thm:generic}
$x_1,\dots,x_n$ を $\mathbb Q$ 上代数的独立な不定元とし，
$e_1,\dots,e_n$ をその基本対称式とする．
\[
 K=\mathbb Q(e_1,\dots,e_n),\qquad
 L=\mathbb Q(x_1,\dots,x_n)
\]
と置く．多項式
\[
 F(T)=\prod_{j=1}^n(T-x_j)
 =T^n-e_1T^{n-1}+e_2T^{n-2}
 -\cdots+(-1)^ne_n
\]
は $K[T]$ に属し，
その分解体は $L$ である．さらに
\[
 \operatorname{Gal}(L/K)\cong S_n.
\]
\end{theorem}
\begin{proof}
$F$ の係数は基本対称式なので $K$ に属する．
根は互いに異なる不定元 $x_1,\dots,x_n$ であり，
標数0だから分離的である．
$L$ はこれらの根によって生成されるので
$F$ の分解体であり，$L/K$ は有限Galoisである．

任意の $\sigma\in S_n$ は
不定元の置換
$x_i\mapsto x_{\sigma(i)}$ によって
$\mathbb Q$ 上の体自己同型 $L\to L$ を定める．
基本対称式は置換で不変なので，
この自己同型は $K$ を固定する．
異なる置換は少なくとも一つの $x_i$ の像が異なるので，
$S_n\hookrightarrow\operatorname{Gal}(L/K)$ は単射である．
逆に任意の $\tau\in\operatorname{Gal}(L/K)$ は
多項式 $F$ の根全体を置換するので，
$\operatorname{Gal}(L/K)\hookrightarrow S_n$ が得られる．
二つの単射から位数は等しく，
群として $\operatorname{Gal}(L/K)\cong S_n$ となる．
\end{proof}

\begin{remark}
$e_1,\dots,e_n$ は代数的独立である．
実際 $L$ は $K$ の有限代数拡大なので
$\operatorname{trdeg}_{\mathbb Q}K
=\operatorname{trdeg}_{\mathbb Q}L=n$ となる．
$K$ は $n$ 個の $e_i$ で生成されるため，
それらの間に非自明な代数関係は存在しない．
従って $F$ は係数が互いに独立な
\emph{一般の $n$ 次多項式}を表す．
\end{remark}

\begin{theorem}[Abel--Ruffiniの定理]\label{thm:abel-ruffini}
$n\ge5$ とする．
一般の $n$ 次方程式
\[
 T^n+a_1T^{n-1}+\cdots+a_n=0
\]
の全ての根を，
係数 $a_1,\dots,a_n$ から四則演算と
有限回の根号抽出によって表す公式は存在しない．
\end{theorem}
\begin{proof}
係数を代数的独立な不定元として扱い，
$K=\mathbb Q(a_1,\dots,a_n)$ を係数体とする．
定理\ref{thm:generic}により，
この一般多項式の分解体のGalois群は $S_n$ である．
$n\ge5$ なら系\ref{cor:sn-nonsolvable}から
$S_n$ は可解群でない．
一方，もし全ての根を根号抽出で表す公式が存在すれば，
多項式は定義によって根号可解であり，
定理\ref{thm:iff}からGalois群は可解でなければならない．
これは矛盾する．
従って一般の方程式に対する根号公式は存在しない．
\end{proof}

\section{3次・4次方程式の根号表示}

\begin{proposition}[3次方程式の根号可解性]\label{prop:cubic}
標数0の体 $K$ 上の任意の3次方程式は根号可解である．
\end{proposition}
\begin{proof}
変数変換によって
$x^3+px+q=0$ の形へ帰着できる．
$p=0$ なら $x^3=-q$ なので根号で解ける．
$p\ne0$ とし $x=u+v$ と置くと
\[
 (u+v)^3+p(u+v)+q
 =u^3+v^3+(3uv+p)(u+v)+q.
\]
そこで $uv=-p/3$，
$u^3+v^3=-q$ とすればよい．
$U=u^3,V=v^3$ は
\[
 U+V=-q,\qquad UV=-p^3/27
\]
を満たすので，$U,V$ は二次方程式
\[
 Z^2+qZ-\frac{p^3}{27}=0
\]
の二根である．
二次方程式の公式によって $U,V$ は平方根で表される．
$U$ の三乗根 $u$ を一つ選び，
$v=-p/(3u)$ と定めれば
$v^3=V$ となり，$x=u+v$ は所望の根である．
三乗根の選択を1の三乗根で変えることにより
三つの根を得られる．
従って分解体は根号拡大の塔に含まれる．
\end{proof}

\begin{proposition}[4次方程式の根号可解性]\label{prop:quartic}
標数0の体 $K$ 上の任意の4次方程式は根号可解である．
\end{proposition}
\begin{proof}
変数の平行移動によって
\[
 x^4+px^2+qx+r=0
\]
の形へ帰着する．
$q=0$ の場合は $z=x^2$ と置くことで
二次方程式に帰着し，根号で解ける．
$q\ne0$ の場合，
\[
 (x^2+ux+v)(x^2-ux+w)
 =x^4+(v+w-u^2)x^2+u(w-v)x+vw
\]
という因数分解を考える．
係数を比較すると
\[
 v+w=p+u^2,\qquad u(w-v)=q,\qquad vw=r.
\]
$z=u^2\ne0$ と置き，
\[
 v=\frac{p+z-q/u}{2},\qquad
 w=\frac{p+z+q/u}{2}
\]
とすれば，$vw=r$ は
\[
 z(p+z)^2-4rz-q^2=0,
\]
すなわち
\[
 z^3+2pz^2+(p^2-4r)z-q^2=0
\]
と同値である．
この3次方程式は命題\ref{prop:cubic}によって根号で解ける．
$q\ne0$ なので $z=0$ は根でなく，
根号表示された非零の根 $z$ を一つ選べる．
$u^2=z$ を満たす $u$ を選び，
上の式で $v,w$ を定めれば，
もとの4次多項式は二つの2次式の積に分解される．
従ってそれぞれ二次方程式の公式で根号表示でき，
全ての根を含む根号拡大の塔が得られる．
\end{proof}

\begin{remark}
Abel--Ruffiniの定理は，
全ての5次以上の方程式が根号可解でないと主張するものではない．
例えば $x^5-2=0$ は
$\sqrt[5]{2}$ と1の五乗根によって解ける．
非可解性は，\emph{一般の}方程式のGalois群が
非可解な対称群となることから導かれる．
\end{remark}

\begin{example}[根号可解な5次方程式]
$f(x)=x^5-2\in\mathbb Q[x]$ とする．
その全ての根は
\[
 \sqrt[5]{2},\quad\zeta_5\sqrt[5]{2},\quad
 \zeta_5^2\sqrt[5]{2},\quad
 \zeta_5^3\sqrt[5]{2},\quad
 \zeta_5^4\sqrt[5]{2}
\]
であり，根号可解である．
\end{example}
\begin{proof}
$\alpha=\sqrt[5]{2}$，$\zeta_5=e^{2\pi i/5}$ とすると
$\alpha^5=2$，$\zeta_5^5=1$ である．
従って表示した各元は $f$ の根であり，
互いに異なるので5次多項式の全ての根を尽くす．
また
\[
 \mathbb Q\subset\mathbb Q(\zeta_5)
 \subset\mathbb Q(\zeta_5,\alpha)
\]
は $\zeta_5^5=1$，$\alpha^5=2$ を満たす根号拡大の塔である．
\end{proof}

\begin{remark}
群論の可解性と体の根号拡大の存在を結びつけるには，
根号の共役を全て加える正規化と，
1の冪根を用いる巡回拡大の構造が必要である．
それらを経由することで，
多項式の根号可解性は分解体のGalois群の
可解性という純代数的な条件へ置き換えられる．
\end{remark}

\section{到達課題の完全解答}
\noindent\textbf{三次多項式のGalois群を決定する．}
\begin{proposition}
$f(x)=x^3-2\in\mathbb Q[x]$ の分解体 $L$ は
$L=\mathbb Q(\sqrt[3]2,\omega)$
（$\omega=e^{2\pi i/3}$）で，
$\operatorname{Gal}(L/\mathbb Q)\cong S_3$ である．
\end{proposition}
\begin{proof}
Eisensteinの判定法から $x^3-2$ は既約なので
$[\mathbb Q(\sqrt[3]2):\mathbb Q]=3$ である．
$\mathbb Q(\sqrt[3]2)\subset\mathbb R$ に
非実数 $\omega$ を添加すると次数2の拡大となる．
従って $[L:\mathbb Q]=6$．
この拡大は標数零の既約多項式の分解体なのでGalois拡大であり，
Galois群は三つの根への忠実な作用から $S_3$ に埋め込まれる．
位数は6なので埋め込みは同型となる．
\end{proof}
\begin{corollary}
$x^3-2=0$ は根号で解ける．
\end{corollary}
\begin{proof}
$S_3$ は正規列
$\{1\}\triangleleft A_3\triangleleft S_3$ を持ち，
商は位数3と2の巡回群なので可解群である．
Galois群が可解であるため根号可解となる．
\end{proof}
\noindent\textbf{到達確認．}
$x^2-2$ のGalois群を求めよ．
\par\noindent\textbf{解答．}
分解体は $\mathbb Q(\sqrt2)$，次数は2なので
Galois群は $\mathbb Z/2\mathbb Z$ に同型である．
\end{document}